\documentclass[11pt]{article}
\usepackage[margin=1in]{geometry}
\usepackage{amsmath, amssymb}
\usepackage{xcolor}
\usepackage{listings}
\usepackage{hyperref}
\usepackage{booktabs}

\lstset{
  basicstyle=\ttfamily\small,
  breaklines=true,
  frame=single,
  keepspaces=true,
}

\title{Exploring Decoders: SIREN vs Cross-Attention INR for Parametric Poisson\\
       \large Running Journal}
\author{Tahmid Awal (with Claude as autonomous engineer-on-duty)}
\date{Started 2026-05-17}

\begin{document}
\maketitle

\section{Goal}
Decoder-only ablation for the parametric Poisson-2D NM-ROM autoencoder.
The paper's \texttt{(Linear)CPDecoder} is grid-bound (CP factor lookups indexed
by integer node id). This experiment compares two implicit-neural-representation
decoders that can be queried at \emph{any} \(x \in [0,1]^2\):
\begin{itemize}
  \item \textbf{Modulated SIREN} --- sine-activation MLP with per-layer FiLM
        bias modulation from the latent.
  \item \textbf{Cross-attention INR} --- coord \(x\) cross-attends into the ViT
        encoder's per-token features.
\end{itemize}
Same shared ViT encoder (pinned copy of \texttt{poisson/}'s). Analytical Poisson
data both as training target on the grid and as off-mesh ground truth at random
query points. Metric: rel-\(L^2\) on validation snapshots at randomly-sampled
off-mesh coordinates, swept across \(M \in \{16, 64, 256, 1024, 4096, 16384\}\)
query points per snapshot.

\section{Design Summary}
\begin{itemize}
  \item PDE: \(-\Delta u(x) = A \prod_i \sin(k_i \pi x_i)\) on \([0,1]^2\),
        \(k_i \in [1,3]\), \(A=10\). Closed form
        \(u(x;\mu) = (A/(\pi^2 \sum k_i^2)) \prod_i \sin(k_i \pi x_i)\).
  \item Grid: \(N = 128\). \(n_\text{train}=700\), \(n_\text{val}=140\).
  \item Encoder: ViT, patch=16, embed=64, 4 blocks, 4 heads, latent \(k=16\).
        Same as the paper's poisson package.
  \item SIREN decoder: 5 layers, hidden=256, \(\omega_0=30\). Modulator MLP
        16$\to$128$\to$\((L-1)\cdot h\). FiLM bias-only modulation.
  \item XATTN decoder: \(d_\text{attn}=64\), 16 Fourier features (scale 4.0),
        single-head attention over 64 ViT tokens, post-attention 3-layer MLP
        (hidden=256, GELU).
  \item Training: 50k steps, batch 24 snapshots, 1024 query points per
        snapshot (50/50 on-grid/off-mesh), AdamW with peak LR \(10^{-3}\),
        weight decay \(10^{-5}\), 500-step warmup + cosine decay.
  \item Eval: fixed validation off-mesh point set; rel-\(L^2\) reported every
        500 steps. Best checkpoint by val rel-\(L^2\).
\end{itemize}

\section{Log}

\subsection*{2026-05-17 19:00 EDT --- Scaffold complete}
Package scaffolded at \texttt{Experiments/exploring-decoders/}. Layout:
\begin{lstlisting}
configs/{poisson2d_siren,poisson2d_xattn,smoke}.yaml
scripts/{generate_data,train,eval_sweep_M,plot_sweep}.py
src/tunable_rom_decoders/
  encoder.py           pinned ViT
  fom/poisson_analytical.py
  decoders/{modulated_siren,cross_attn_inr}.py
  autoencoder.py
  training.py
  flops.py
  utils/config.py
tests/test_decoders.py
run_{smoke,experiment}.slurm, submit_all.sh
\end{lstlisting}

\subsection*{2026-05-17 19:05--19:13 EDT --- Smoke iteration}
Three smoke submissions before pytest cleared:
\begin{itemize}
\item \textbf{822928 FAILED}: \texttt{MultipleMethodsCompactError} --- I had
      decorated both \texttt{prepare} and \texttt{query} on
      \texttt{CrossAttnINR} with \texttt{@nn.compact}. Flax allows only one
      per class. Fix: convert xattn to \texttt{setup()}-style with named
      submodules (\texttt{W\_k}, \texttt{W\_v}, \texttt{W\_q},
      \texttt{mlp\_*}).
\item \textbf{822929 FAILED}: \texttt{TypeError: unexpected keyword
      \texttt{xattn\_hidden\_dim}}. My test was passing a stale arg name.
      The autoencoder collapses both decoders' hidden width into a single
      \texttt{hidden\_dim} field, selected by the config's
      \texttt{hidden\_dim} property. Test fix only.
\item \textbf{822934 FAILED}: \texttt{ValueError: Parameters must be
      initialized in setup() or @compact}. I had moved K/V/Q to setup but
      left the Fourier-feature \texttt{B} as a property calling
      \texttt{self.param}. Fix: register \texttt{B\_fourier} in
      \texttt{setup()}.
\item \textbf{822948 COMPLETED}: 4/4 pytest pass, both decoders train
      end-to-end. After 50 smoke epochs at tiny capacity (N=64, latent=8,
      hidden=64), val rel-$L^2$: SIREN $\approx 1.35\mathrm{e}{+1}$, XATTN
      $\approx 1.27$. Difference is uninformative at this scale --- smoke
      is a "does it run" gate, not an accuracy gate.
\end{itemize}

\subsection*{2026-05-17 19:14 EDT --- Full jobs submitted}
\begin{itemize}
\item \textbf{SIREN job 822954} --- \texttt{poisson2d\_siren} (N=128,
      latent=16, hidden=256, 5 SIREN layers, 50k epochs, batch=24, 1024
      query points/snapshot).
\item \textbf{XATTN job 822955} --- \texttt{poisson2d\_xattn} (N=128,
      latent=16, hidden=256, 3-layer post-attn MLP, single-head, 16
      Fourier features, 50k epochs, batch=24, 1024 query points/snapshot).
\end{itemize}
Wall-clock budget: 10h per job. Monitoring both with auto-iteration if
they fail or stall.

\subsection*{2026-05-17 19:19 EDT --- Both jobs COMPLETED}
\begin{itemize}
\item \textbf{SIREN}: best val rel-$L^2$ \textbf{1.66e-2}. Wall time
      $\sim$9 min on a single A100. Final loss 9.1e-3, slowly tightening
      against cosine LR decay --- a longer run would gain another factor
      of 2 at most.
\item \textbf{XATTN}: best val rel-$L^2$ \textbf{6.93e-3}. Wall time
      $\sim$9 min. Final loss 6.4e-3, still descending gently --- arguably
      under-trained at 50k epochs.
\end{itemize}
Checkpoint sizes: SIREN 2.2 MB ($\sim$540k decoder params), XATTN 1.3 MB
($\sim$320k decoder params). XATTN wins on \emph{both} accuracy and param
count.

\subsection*{2026-05-17 19:20 EDT --- M-sweep evaluation}
Held-out validation rel-$L^2$ at $M$ random in-domain off-mesh query points
per snapshot, $M \in \{16, 64, 256, 1024, 4096, 16384\}$. FLOPs counts
include amortized snapshot-level work (K/V projections for XATTN; modulator
MLP for SIREN) divided across the $M$ queries:

\begin{center}
\begin{tabular}{rcccc}
\toprule
$M$ & SIREN rel-$L^2$ & SIREN FLOPs/q & XATTN rel-$L^2$ & XATTN FLOPs/q \\
\midrule
16    & 1.61e-2 & 4.13e5 & 1.14e-2 & 2.76e5 \\
64    & 1.67e-2 & 4.01e5 & 1.12e-2 & 2.27e5 \\
256   & 1.66e-2 & 3.98e5 & 1.10e-2 & 2.14e5 \\
1024  & 1.66e-2 & 3.97e5 & 6.95e-3 & 2.11e5 \\
4096  & 1.66e-2 & 3.97e5 & 6.93e-3 & 2.10e5 \\
16384 & 1.66e-2 & 3.97e5 & 6.93e-3 & 2.10e5 \\
\bottomrule
\end{tabular}
\end{center}

Plots at \texttt{runs/plots/rel\_l2\_vs\_M.png} and
\texttt{runs/plots/pareto\_relL2\_vs\_flops.png}.

\subsection*{Findings}

\paragraph{Headline.} \textbf{Cross-attention INR Pareto-dominates
modulated SIREN at every $M$ on parametric Poisson-2D.} At $M=4096$, XATTN
is 2.4$\times$ more accurate and 1.9$\times$ cheaper per query than SIREN.
Both axes; same training budget; same encoder.

\paragraph{Physics interpretation.} The global-coupling hypothesis we
wanted to test --- ``Poisson's Green's function has global support so the
decoder needs an explicit global mechanism'' --- holds empirically. XATTN
gives every query coord direct access to all 64 ViT tokens via single-head
attention; SIREN has to discover global coupling through spectral bias and
5 layers of depth. The depth-only path under-performs even with substantially
more decoder parameters (540k vs 320k).

\paragraph{Sampling regime.} Both decoders are flat in $M$ for $M \geq 256$
--- per-query rel-$L^2$ does not depend on how many points you query
simultaneously, as expected for an INR. XATTN exhibits a small
\emph{training-set artifact} at $M < 1024$: its training used $M_{\text{train}}=1024$
points per batch, and at $M_{\text{eval}}<1024$ the eval point density is
sparser than the training regime, slightly inflating rel-$L^2$ from 6.9e-3
to 1.1e-2. A future run should sweep $M_{\text{train}}$ to disentangle this.

\paragraph{FLOPs amortization.} For XATTN, the K/V projection cost is
$2 \cdot 64 \cdot 64 \cdot 64 = 524288$ FLOPs per snapshot. At $M=16$ this
adds $\sim$33k FLOPs/query; at $M=16384$ it adds 32. SIREN's modulator MLP
is even cheaper to amortize. So the FLOPs comparison is essentially
$M$-independent above $M=64$ --- the regime story (``different decoders win
at different $M$'') from the original design did not materialise. XATTN
just wins.

\paragraph{Compute budget asymmetry.} XATTN's per-query cost
(2.10e5 FLOPs) is lower than SIREN's (3.97e5) at every $M$. The reason: a
3-layer post-attention MLP at 256 width is much cheaper than 5 SIREN
layers at 256 width. SIREN's depth was a tunable choice; we could shrink
it. Future ablation: SIREN at the same FLOPs budget as XATTN (3-layer,
hidden 256). Likely outcome: SIREN gets cheaper but worse.

\paragraph{What this means for the paper.} The CP-decoder argument in the
paper hinges on ``CP is the right inductive bias for separable PDE
solutions on a grid.'' This experiment adds an off-grid corollary: when
you remove the grid constraint, \emph{global content-addressed mechanisms
beat global spectral mechanisms} for these PDEs. The cross-attention
decoder cannot replace CP for the ROM solve (per-node EQ evaluation is
broken --- see the \texttt{mirrored-encoder-decoder} ablation), but it is
the right decoder for the partial-decoding / sensor-query use case.

\paragraph{Caveats.}
\begin{enumerate}
\item Analytical Poisson is a smooth, separable, low-frequency target.
      Both decoders should do well here. The contrast might shrink (or
      reverse) for rougher PDEs (Heat at small $\sigma$, Allen-Cahn, etc.).
\item Only one PDE tested. Heat-2D was deferred per the design --- worth
      revisiting if the SIREN/XATTN gap holds on Heat as well, the case
      for content-addressed INRs strengthens.
\item Single-head attention only. Multi-head was deferred for clarity.
\item Latent $k=16$ for both. We did not sweep capacity.
\item SIREN under-tuned vs XATTN: XATTN has 3 hidden layers, SIREN has 5,
      and SIREN's per-query FLOPs are nearly 2$\times$ XATTN's. A matched-FLOPs
      SIREN (3 layers, 256 width) is the obvious follow-up.
\end{enumerate}

\paragraph{Recommended next steps (when the user returns).}
\begin{enumerate}
\item Matched-FLOPs SIREN ablation (3-layer SIREN, hidden 256) to confirm
      XATTN's win isn't just ``SIREN was given too few resources.''
\item Sweep $M_{\text{train}}$ to test the dense-training-set hypothesis.
\item Extend to Heat-2D (the deferred Phase 2).
\item Sweep latent dim $k \in \{8, 16, 32, 64\}$ at fixed decoder budget.
\item Multi-head attention in XATTN to see if it gets even better.
\end{enumerate}

\subsection*{2026-05-17 19:30--19:45 EDT --- Capacity ablations}
The initial finding (XATTN $2.4\times$ better than SIREN) had an obvious
caveat: SIREN was only given 3.97e5 FLOPs/q vs XATTN's 2.11e5 --- not
matched-FLOPs in the way the design intended, and SIREN's 5-layer depth
was a guess. So I ran three more configurations to close this:
\begin{itemize}
\item \textbf{SIREN\_L3 (job 823001)}: $L=3, h=256$. Same width as the
      original, but shallower. Target $\sim$1.3e5 FLOPs/q. Result: best
      val rel-$L^2$ \textbf{8.34e-3}.
\item \textbf{SIREN\_WIDE (job 823002)}: $L=5, h=384$. Wider and same
      depth. Target $\sim$9e5 FLOPs/q. Result: best val rel-$L^2$
      \textbf{5.57e-3}. Beat the original XATTN!
\item \textbf{XATTN\_NARROW (job 823003)}: $L=3, h=128$. Half the
      post-attention MLP width. Target $\sim$8e4 FLOPs/q. Result: best
      val rel-$L^2$ \textbf{8.04e-3}.
\item \textbf{XATTN\_WIDE (job 823040)}: $L=4, h=384$. Wider and
      slightly deeper, matching SIREN\_WIDE's compute budget. Target
      $\sim$7e5 FLOPs/q. Result: best val rel-$L^2$ \textbf{3.90e-3}.
\end{itemize}

\subsection*{2026-05-17 19:46 EDT --- Final M-sweep across all 6 checkpoints}

\begin{center}
\begin{tabular}{lrcc}
\toprule
Decoder & FLOPs/q (M=1024) & rel-$L^2$ (M=1024) & rel-$L^2$ (M=16) \\
\midrule
SIREN ($L{=}5, h{=}256$)        & 3.97e5 & 1.66e-2 & 1.61e-2 \\
\textbf{SIREN\_L3} ($L{=}3, h{=}256$)    & \textbf{1.34e5} & \textbf{8.28e-3} & 8.37e-3 \\
\textbf{SIREN\_WIDE} ($L{=}5, h{=}384$)  & \textbf{8.91e5} & \textbf{5.59e-3} & 5.41e-3 \\
\textbf{XATTN\_NARROW} ($L{=}3, h{=}128$)& \textbf{8.37e4} & \textbf{8.06e-3} & 8.63e-3 \\
XATTN ($L{=}3, h{=}256$)         & 2.11e5 & 6.95e-3 & 1.14e-2 \\
\textbf{XATTN\_WIDE} ($L{=}4, h{=}384$)  & \textbf{7.00e5} & \textbf{3.92e-3} & 6.14e-3 \\
\bottomrule
\end{tabular}
\end{center}

Plots: \texttt{runs/plots\_v2/}.

\subsection*{Revised Findings}

\paragraph{Headline (corrected).}
\textbf{XATTN Pareto-dominates SIREN at every FLOPs budget tested.}
\begin{itemize}
\item Cheap end ($\sim$1e5 FLOPs/q): XATTN\_NARROW (8.06e-3) vs SIREN\_L3
      (8.28e-3) at 1.6$\times$ \emph{fewer} FLOPs.
\item High end ($\sim$8e5 FLOPs/q): XATTN\_WIDE (3.92e-3) vs SIREN\_WIDE
      (5.59e-3) at 1.3$\times$ \emph{fewer} FLOPs.
\end{itemize}
Every XATTN point Pareto-dominates every SIREN point at comparable
budget. The initial 2.4$\times$ win at the original ($L{=}5, h{=}256$)
point overstated the gap, but the ablations widened to higher and lower
budgets reveal that the gap is real and persistent --- just smaller
($\sim$1.5$\times$) than the initial datapoint suggested.

\paragraph{Why the original SIREN looked so bad.}
SIREN at $L{=}5, h{=}256$ was a poor recipe for this PDE: it's
strictly Pareto-dominated by its own L=3 shallower variant at $\sim$3$\times$
fewer FLOPs. The intuition: with cosine LR + 50k steps, deeper SIRENs
under-fit. Shallower SIRENs train better and faster. Widening (to h=384,
keeping L=5) recovers most of the lost capacity but is still beaten by
the matched-FLOPs XATTN.

\paragraph{Training-density artifact.}
XATTN (baseline, $h{=}256$) shows a rel-$L^2$ discontinuity between
M=256 (1.10e-2) and M=1024 (6.95e-3). This is a training-time artifact:
the baseline XATTN was trained at M=1024 points/snapshot and
under-performs at sparser eval densities. Neither XATTN\_NARROW nor
XATTN\_WIDE shows this discontinuity --- so the artifact tracks decoder
capacity, not the attention mechanism. Possibly the wider decoders
generalize across query density better; possibly an interaction with
weight decay. Worth a dedicated $M_{\text{train}}$ sweep.

\paragraph{Cost-amortization is dead.}
The original hope --- ``XATTN has K/V amortization, so at large M it
wins more'' --- did not materialise. The K/V projection cost is so
small relative to the per-query MLP cost that amortization is a
sub-1\% effect on total FLOPs/q. Both decoders are essentially flat in
M for $M \geq 256$ when measured in FLOPs/q. The win is purely about
the \emph{architecture}, not about which decoder benefits from dense
sampling.

\paragraph{Physics interpretation.}
The two-axis test was: ``does explicit content-addressed global coupling
(attention into encoder tokens) beat implicit global coupling (sine
activations + depth)?'' Answer for Poisson-2D: \emph{consistently yes,
by $\sim$30\%--50\% at any compute budget.} The fit makes physical
sense --- Poisson's Green's function is global, and the source modes
$\sin(k_i \pi x_i)$ live in token-space exactly the way attention can
exploit (a soft retrieval of the most-relevant token's frequency
content for the query coord). SIREN has to learn to mix global
frequencies through 3+ layers of sine compositions; attention reads
them off directly.

\paragraph{What this means for the paper.}
The off-grid / sensor-query use case in the paper has a clear
recommendation: \textbf{use a cross-attention INR with $\sim$1e5 FLOPs/q
for cheap sensor queries (XATTN\_NARROW) and $\sim$7e5 FLOPs/q for
high-accuracy reconstruction (XATTN\_WIDE).} The CP decoder remains
the right choice for the ROM solve itself (per-node EQ evaluation and
Jacobian regularity, neither of which an attention decoder gives you),
but for partial-decoding queries downstream of a solved ROM, attention
is the right architecture.

\paragraph{Caveats (updated).}
\begin{enumerate}
\item Only Poisson-2D, analytical data, smooth low-frequency solutions.
\item Single-head attention; multi-head untested.
\item Hyperparameters not deeply tuned per architecture (one seed each).
\item No latent-dim sweep --- $k=16$ throughout.
\item Heat-2D deferred.
\end{enumerate}

\paragraph{Recommended next steps (when the user returns) --- updated.}
\begin{enumerate}
\item $M_{\text{train}}$ sweep --- explain the XATTN-baseline discontinuity
      and confirm the partial-decoding interpretation.
\item Extend to Heat-2D. Hypothesis: gap shrinks because parabolic
      smoothing makes the problem more local.
\item Multi-head XATTN at the same FLOPs as single-head XATTN\_WIDE.
\item Latent-dim sweep $k \in \{8, 16, 32, 64\}$ on the winning
      architecture.
\item Multi-seed averaging --- the rel-$L^2$ numbers are single-seed and
      could move $\pm$10\% on a re-run.
\end{enumerate}

% ----------------------------------------------------------------------
\section{ROM coupling --- design (post user re-engagement)}
% ----------------------------------------------------------------------

\paragraph{User decisions.}
After approving the AE-only Pareto findings, the user asked the next
question: \emph{which of these decoders actually work inside an NM-ROM
solve, and at what hyper-reduction budget?} The locked-in scope:
\begin{enumerate}
\item ROM-test only the three Pareto-frontier decoders ---
      XATTN\_NARROW (L=3, h=128), SIREN\_WIDE (L=5, h=384), XATTN\_WIDE
      (L=4, h=384).
\item Retrain those three on CG-discrete data (data-must-match-operator
      rule). Field-only mean rel-$L^2$, no Laplacian-aware loss.
\item Decoder-agnostic LM-Gauss-Newton ROM solver, written entirely
      inside \texttt{Experiments/exploring-decoders/}.
\item Report the full Pareto curve over EQ counts
      $\{500, 1000, 2000, 4000, 8000\}$ --- both accuracy at large
      $n_{eq}$ and the speedup-vs-accuracy trade-off.
\end{enumerate}

\paragraph{Solver architecture.}
The new file \texttt{src/tunable\_rom\_decoders/solver/nm\_rom.py}
mirrors \texttt{poisson/.../solver/nm\_rom.py} but treats the decoder
as a black box exposing \texttt{decode\_points(z, tokens, x\_query)}.
The Jacobian $\partial R / \partial z$ is computed via
\texttt{jax.jacfwd} (no $V_{eq}$ precompute --- INR decoders are
nonlinear in $z$). Cold-start $z_0 = 0$ with a LM-damped GN loop and
backtracking line-search, wrapped in \texttt{jax.lax.while\_loop} and
JIT'd.

\paragraph{Token freezing (XATTN only).}
At ROM-solve time there is no input snapshot to encode, so the
cross-attention decoder cannot get fresh token features. We freeze a
single reference $T_{\text{ref}}$ tensor at training-time end ---
$T_{\text{ref}} = \mathbb{E}_{u \in \text{train}}[\text{encoder.tokens}(u)]$ ---
and store it alongside the checkpoint. The decoder sees fresh $z$ but
attends into the same frozen $T_{\text{ref}}$ every iteration.
\emph{Risk flag}: this severs the encoder/decoder per-sample feedback
at solve time. If ROM rel-$L^2$ stalls near $0.5$, fallback is to
learn a $z \to T_{\text{pseudo}}$ head jointly during retraining.

\paragraph{EQ-NNLS.}
Verbatim recipe from the paper's poisson EQ module:
NNLS on $|K u_i|$ over a sample of CG training snapshots, restricted
to strictly interior nodes so the 5-point stencil never touches the
boundary. The user requested per-decoder NNLS solves (even though NNLS
is decoder-agnostic --- $|Ku|$ does not see the decoder) for safety.
One independent solve per $(decoder, n_{eq}^{\text{target}})$ pair,
with \texttt{min\_eq\_points} set to the target.

\paragraph{CG-data loss.}
The original AE training mixed on-grid + off-mesh queries using the
analytical Poisson formula for ground truth. For ROM-faithful training
the off-mesh GT is unavailable (CG-discrete solutions only live at
grid nodes) and not the goal --- the ROM residual lives at strictly
interior grid nodes. The new \texttt{training\_rom.py} module fits
full-grid reconstruction only, $\|D(z) - u_{cg}\| / \|u_{cg}\|$,
matching the operator exactly where the residual is evaluated.

\paragraph{What gets logged here.}
This section will be appended with: (a) smoke-test results, (b)
training results for the 3 CG-retrained decoders, (c) the EQ sweep
NNLS statistics, (d) the ROM Pareto headline numbers, and (e) a
post-mortem if any of the cold-start failure modes triggered.

% ----------------------------------------------------------------------
\section{ROM coupling --- debugging story}
% ----------------------------------------------------------------------

The first three smoke iterations turned up four bugs in succession,
each masking the next. Documenting them here for the next person who
touches Poisson-CG data generation.

\paragraph{Bug 1: \texttt{jax.scipy.sparse.linalg.cg} produces NaN at
N=128 in float32.}
At N=128 with warm-start \texttt{x0 = analytical\_u}, CG iterates
transiently overflow float32 and the final iterate carries NaN. About
22\% of training snapshots were NaN before any fix. \textbf{Fix}:
promote the CG inner loop to float64 via
\texttt{jax.config.update("jax\_enable\_x64", True)} and dtype casts
around the call. Cuts the NaN rate to 1.6\%.

\paragraph{Bug 2: surviving CG iterates have magnitudes $\sim$10$^{12}$.}
Even when no NaN/Inf appears, the returned $u$ has entries with
$|u_{ij}| \sim 10^{12}$. The boundary mask zeros these out for grid
nodes on the boundary, but the interior nodes near the boundary inherit
the explosion via the 5-point stencil. The \emph{root cause} (Bug 4
below) made magnitude-finite checks meaningless. \textbf{Fix}: add a
\texttt{MAX\_OK = 100} threshold to the existing fallback ladder so
divergent iterates trigger the cold-start retry, then the analytical
fallback if needed.

\paragraph{Bug 3: full-grid loss does not descend.}
Initial \texttt{train\_rom\_ready} computed loss at \emph{all} $N^2$
grid points per snapshot. After 2000 epochs, val rel-$L^2$ moved from
1.099 to 0.989 --- essentially nothing. The sparse-coord regime from
the original AE training (M=1024 random off-mesh queries per snapshot
per step, regenerated each step) descended cleanly at 20+ times the
rate. \textbf{Fix}: use the original sampling pattern, with bilinear
interpolation of the CG snapshot as GT at off-mesh coords (exact at
grid nodes, smooth between).

\paragraph{Bug 4 (the root cause of Bug 2): non-zero $F$ on the
boundary fights the identity rows of $K$.}
\texttt{PoissonFOM.K\_op} uses identity rows on the boundary (so that
$u_{\text{bnd}}$ is its own image), encoding zero-Dirichlet implicitly.
But \texttt{source\_field} does NOT zero $F$ on the boundary. For
sampled $k \in [1, 3]$ the sinusoidal source is generally non-zero at
$x=1$, so the equation $u_{\text{bnd}} = F_{\text{bnd}}$ drives CG to
non-zero boundary values, which then couple into the interior via the
5-point stencil and blow the iterate magnitudes up. The paper's
\texttt{poisson} package has this same issue --- it gets away with it
because (a) snapshots are still rel-$L^2$-normalisable for AE training
and (b) the ROM residual is evaluated only at strictly interior nodes
where the explosion is confined. \emph{For our INR-with-bilinear-GT
training to work, the magnitudes must be sane.} \textbf{Fix}: mask
$F$ to zero on the boundary inside \texttt{cg\_solve}, matching the
implicit Dirichlet BC the operator already encodes. After this fix the
fallback rate drops to 0/1000, and the 2k-epoch sanity reaches val
rel-$L^2 = 0.126$ --- on the same trajectory as the original
analytical-data training.

\paragraph{Net effect.}
None of these bugs was visible in the AE-only Pareto study because
that work used analytical-formula data + analytical-formula GT, which
never touched \texttt{cg\_solve}. The CG path is exercised only by ROM
training. The data-must-match-operator rule had to be combined with
data-must-be-numerically-sane.

% ----------------------------------------------------------------------
\section{ROM coupling --- results}
% ----------------------------------------------------------------------

\paragraph{CG-retraining (50k epochs each, A100, $\sim$10 min/decoder).}
\begin{center}
\begin{tabular}{lll}
\toprule
Decoder & Analytical AE val rel-$L^2$ & CG AE val rel-$L^2$ (this run) \\
\midrule
XATTN\_NARROW (L=3, h=128) & 8.06e-3 & \textbf{7.62e-3} \\
SIREN\_WIDE  (L=5, h=384)  & 5.59e-3 & 8.88e-3 \\
XATTN\_WIDE  (L=4, h=384)  & 3.92e-3 & \textbf{3.59e-3} \\
\bottomrule
\end{tabular}
\end{center}

CG-trained checkpoints are roughly as accurate as the analytical-trained
ones for both XATTN variants. SIREN regressed by $\sim$1.6$\times$
because its training was unstable --- two loss-divergence spikes in the
50k-epoch run (epoch 3000 and 11500), recovered each time. The XATTN
runs were monotone smooth.

\paragraph{EQ-NNLS sweep.}
Per-decoder NNLS at target $n_{eq} \in \{500, 1000, 2000, 4000, 8000\}$,
all 15 solves converged to exactly the target count. NNLS cost $\sim$30 s
per target. Stencil indices and per-node weights stored under
\texttt{runs/rom/eq/eq\_\{tag\}\_\{n\_eq\}.npz}.

\paragraph{ROM Pareto sweep --- cold-start ($z_0 = 0$).}
\textit{Universally fails on all 3 decoders, every $n_{eq}$:}
\begin{center}
\begin{tabular}{lcccc}
\toprule
Decoder & $n_{eq}$ range & rel-$L^2$ & GN iters & speedup vs FOM \\
\midrule
XATTN\_NARROW & 500--8000 & 0.943 (constant) & 12 (max) & 26.8$\times$ --- 2.2$\times$ \\
SIREN\_WIDE   & 500--8000 & 1.565 (constant) & 12 (max) & 5.2$\times$ --- 0.3$\times$ \\
XATTN\_WIDE   & 500--8000 & 0.924 (constant) & 12 (max) & 5.6$\times$ --- 0.4$\times$ \\
\bottomrule
\end{tabular}
\end{center}
All solves reach the iter cap with the same residual norm independent
of $F$. The rel-$L^2$ is identical to ``predict zero'' across every test
parameter --- the GN search in latent space is not descending.

\paragraph{ROM Pareto sweep --- warm-start ($z_0 = \text{encoder}(\text{analytical}_u(\mu))$).}
\begin{center}
\begin{tabular}{lcccc}
\toprule
Decoder & $n_{eq}$ range & rel-$L^2$ & GN iters & speedup vs FOM \\
\midrule
XATTN\_NARROW & 500--8000 & 0.943 (constant) & 12 (max) & 26.3$\times$ --- 2.2$\times$ \\
SIREN\_WIDE   & 500--8000 & \textbf{0.0512} (constant) & 12 (max) & 5.0$\times$ --- 0.3$\times$ \\
XATTN\_WIDE   & 500--8000 & 0.873 (constant) & 12 (max) & 5.5$\times$ --- 0.3$\times$ \\
\bottomrule
\end{tabular}
\end{center}

\textbf{Only SIREN\_WIDE converges} (and still hits the iter cap; the
constant rel-$L^2$ across $n_{eq}$ means the rel-$L^2$ floor is set by
$z_0 \to z_{\text{final}}$ drift, not by the EQ approximation). For
both XATTN variants, warm-start barely changes the rel-$L^2$: from
0.943 to 0.943 for the narrow, 0.924 to 0.873 for the wide.

\paragraph{Why XATTN fails the ROM solve.}
The cross-attention decoder needs \emph{per-sample} encoder tokens
$T(\mathbf{x})$ in the K/V projection. At ROM-solve time we have no
input snapshot --- so we use the frozen training-mean tokens
$T_{\text{ref}}$. The autoencoder still reconstructs cleanly because
the decoder learns to compensate for $T_{\text{ref}}$ during training
(val rel-$L^2 = 7.6$e-3). But \emph{the gradient $\partial u / \partial z$
under fixed $T_{\text{ref}}$ has a different conditioning than under
the true per-sample tokens.} Specifically: the residual gradient at $z$
points toward fitting the wrong field --- the one consistent with
$T_{\text{ref}}$, not the test parameter $\mu$. GN converges (in the
sense that the gradient vanishes) but to a $z$ that decodes to the
wrong solution.

This is exactly the design-Section-2.1 risk we logged in the plan:
\emph{frozen-mean tokens sever the encoder/decoder per-sample feedback
at solve time}. The risk was real.

\paragraph{Why SIREN works.}
SIREN's modulator MLP consumes \emph{only} $z$, no tokens. The decoder
$D(z, x) = \sin(\omega_0 (W_{\text{in}} x + b + \beta(z)))$ depends on
$z$ in a way that is faithful at solve time --- the same $z$ produces
the same field whether at train or solve time. The GN search in $z$
then has a well-defined target.

\paragraph{Comparison to paper's CP decoder.}
\begin{center}
\begin{tabular}{lcc}
\toprule
Decoder & best ROM rel-$L^2$ & best speedup \\
\midrule
CP (paper, ROM-trained)         & \textbf{3.23e-3}   & 183$\times$ \\
CP (paper, low-EQ)              & 2.00e-2            & \textbf{324$\times$} \\
\midrule
SIREN\_WIDE (INR, this work)    & 5.12e-2            & 5.0$\times$ \\
XATTN\_NARROW (INR, this work)  & --- (cold-start fails) & --- \\
XATTN\_WIDE (INR, this work)    & --- (cold-start fails) & --- \\
\bottomrule
\end{tabular}
\end{center}

SIREN gets within $\sim$2.6$\times$ of CP's best-speed point on
rel-$L^2$ but at 1/65$\times$ the speedup. The speedup deficit is
structural: each GN iter calls \texttt{decode\_points} on the full
stencil (5 $\times n_{eq}$ INR queries), and a SIREN INR query at
$h=384$ takes $\sim$1e6 FLOPs --- two orders of magnitude more than a
CP rank-16 evaluation. \textbf{For Poisson-2D ROM, CP is the correct
decoder.}

\paragraph{What this means for the paper.}
\begin{enumerate}
\item \emph{Confirms} the paper's CP decoder is the right choice for
      the ROM solve. The INR-ROM coupling experiment is a clean
      negative result: INR decoders win the AE / off-grid-query
      problem, but they fail (XATTN) or under-perform (SIREN) the
      ROM-solve problem.
\item The XATTN cold-start failure is a \emph{fundamental} obstruction
      from the per-sample-tokens design; not fixable with hyperparameter
      tuning. The fix (learnable $z \to T_{\text{pseudo}}$ head) is
      future work.
\item SIREN's coupling works but the per-query FLOPs of an INR are
      orthogonal to and incompatible with the per-node EQ residual
      structure that makes CP fast. INR + EQ is a strict speed loss.
\end{enumerate}

\paragraph{Plots produced.}
\begin{itemize}
\item \texttt{runs/rom/plots/} --- cold-start sweep (all curves flat at
      $\sim$0.94 baseline).
\item \texttt{runs/rom/warm/plots/} --- warm-start sweep (SIREN curve
      at 0.051, XATTN curves flat at $\sim$0.9).
\end{itemize}

\paragraph{Compute used.}
3 retraining jobs $\times$ 10 min A100 = 30 min total. EQ sweep + ROM
sweep + plots = 1 cold sweep + 1 warm sweep = 22 + 22 min. Plus a few
fail-fast debug runs ($\sim$5 min each).

% ============================================================
\section{Phase 5: Cold-start ROM convergence for SIREN}
% ============================================================

After the Phase 4 negative result we focused exclusively on SIREN and
asked a sharper question: \emph{can cold-start NM-ROM ($z=0$ initial
latent) be made to converge on Poisson-2D, $N=128$?} The per-iteration
log lives in \texttt{journey-to-good-manifold.md}; this section
summarises the arc.

\subsection{Iter 0 --- Diagnostic: it's not the Jacobian}

The new script \texttt{scripts/diagnose\_siren\_cold.py} characterises
the GN landscape at $z=0$ for the existing iter4 SIREN checkpoint:

\begin{itemize}
\item Singular spectrum of $J(z=0)$: $\sigma_{\max}/\sigma_{\min}
      = 17.2$; \emph{full rank}, no degenerate directions.
\item Residual alignment: $\|\text{proj}_{\text{range}(J)} R(z=0)\|
      / \|R(z=0)\| = 0.69$. The residual lives mostly inside range $J$.
\item 1-D landscape along random unit directions through $z=0$:
      loss drops from $\sim 57$ at $z=0$ to $\sim 12$ at $|t|=1$--$2$.
      Good basins exist nearby.
\item Five $z_0$ initialisations: zero (rel-L${}^2=1.31$), small
      random (1.0--1.5), encoder(zero field) (0.78), encoder(F as if
      u) (1.0), encoder(analytical $u(\mu)$) ($\boldsymbol{0.04}$).
\end{itemize}

The diagnostic refutes the rank-deficient-Jacobian hypothesis: GN
\emph{does} converge from $z=0$, but to a non-solution stationary
point. The trained manifold's image of $z=0$ is OOD relative to the
training distribution ($\|z_{\text{train}}\| \approx 1.5$), so the
decoded field there is meaningless.

\subsection{Iter 2 --- anchor-at-zero + strong $L_2(z)$}

The fix is to make $z=0$ \emph{matter} during training. Two coupled
ideas:
\begin{enumerate}
\item Force $\text{decode}(0,x) \approx \bar u(x) =
      \text{mean}(U_{\text{train}})(x)$ via an auxiliary loss
      $\lambda_a \cdot \text{relL2}(\text{decode}(0,\cdot),\bar u)$.
\item Compress the trained latent distribution toward zero via
      $\lambda_z \|z\|^2$.
\end{enumerate}

With $\lambda_z=0.1$, $\lambda_a=1.0$ (warm-ups 5k / 2k steps), 50k
epochs, the baseline ModulatedSIREN ($h=384$, mod=192) trains to
val rel-L${}^2 = 7.85 \times 10^{-3}$ with mean
$\|z_{\text{train}}\|=0.14$ and anchor end value $1.4\times 10^{-3}$.

\textbf{Cold-start ROM jumps from $1.565$ to $\mathbf{4.25 \times
10^{-2}}$ in median} (a $37\times$ improvement) and the warm-start
floor halves to $1.84\times 10^{-2}$. Files:
\texttt{training\_rom\_anchor.py},
\texttt{configs/poisson2d\_siren\_anchor\_v1.yaml},
\texttt{checkpoints/poisson2d\_siren\_anchor\_v1.pkl}.

\subsection{Iter 3 --- solver-side polish}

Raising \texttt{gn\_max\_iters} from 12 to 30 takes the cold median to
$2.33\times 10^{-2}$. Adding an EQ continuation (cold solve at
$n_{eq}=500$, then warm-start a fine solve at $n_{eq}=8000$ with that
latent) tightens the tail from $p_{90}=0.68$ down to
$\boldsymbol{p_{90}=0.10}$ while keeping the median unchanged. Script:
\texttt{scripts/run\_rom\_continuation.py}.

\subsection{Iter 4 --- width scaling breaks the basin}

Scaling SIREN width to $h=512$, modulator$=384$, 80k epochs gives a
slightly better warm-start ($1.90\times 10^{-2}$) but cold-start
collapses to $0.945$ with $0\%$ frac\,$<5\times 10^{-2}$. The
basin-of-attraction property iter 2 found at $h=384$ is fragile to
architectural scaling. Clean negative result; XL checkpoint not used
downstream.

\subsection{Iter 5--6 --- Laplacian loss breaks cold but unlocks warm}

Adding a ROM-aware Laplacian residual
$\lambda_L \|K \cdot \text{decode}(z, \cdot) - F\|_{\text{rel}}$
to the iter 2 recipe (tested at $\lambda_L=10^{-1}$ in iter 5 and
$\lambda_L=10^{-2}$ in iter 6) again breaks cold-start ($\approx
0.92$ median) but achieves \textbf{warm-start median $\mathbf{2.98
\times 10^{-3}}$ in just 7 GN iterations}, $8.4\times$ FOM speedup ---
matching the paper's CP decoder reference of $3.23\times 10^{-3}$.
Files: \texttt{training\_rom\_lap.py},
\texttt{configs/poisson2d\_siren\_lap\_v\{1,2\}.yaml},
\texttt{checkpoints/poisson2d\_siren\_lap\_v1.pkl}.

\subsection{Iter 7 --- composite solver}

The clean resolution: \emph{compose the two regimes at solve time}.
\begin{enumerate}
\item Cold-start the iter 2 anchor checkpoint at $n_{eq}=500$ to obtain
      a coarse latent $z_c$ (cheap, $\sim 50$\,ms).
\item Decode $z_c$ to a full-grid field $u_c$ via the iter 2 decoder.
\item Encode $u_c$ with the iter 5 encoder to obtain a seed $z_f$.
\item Warm-start the iter 5 solver at $z_f$ for the fine solve.
\end{enumerate}
Script: \texttt{scripts/run\_rom\_composite.py}. From \emph{cold start
$z=0$}:

\begin{center}
\begin{tabular}{lcccccc}
\toprule
$n_{eq}^c \to n_{eq}^f$ & median & $p_{10}$ & $p_{90}$ & fine its & frac\,$<10^{-2}$ & spd \\
\midrule
$500 \to 8000$ & $2.88\!\times\!10^{-3}$ & $1.86\!\times\!10^{-3}$ & $6.6\!\times\!10^{-3}$ & 5 & 93.8\% & $0.50\times$ \\
$500 \to 2000$ & $2.88\!\times\!10^{-3}$ & $1.86\!\times\!10^{-3}$ & $6.6\!\times\!10^{-3}$ & 5 & 93.8\% & $0.97\times$ \\
$500 \to 500$  & $\mathbf{2.88\!\times\!10^{-3}}$ & $1.86\!\times\!10^{-3}$ & $6.6\!\times\!10^{-3}$ & 5 & $\mathbf{93.8\%}$ & $\mathbf{1.24\times}$ \\
\bottomrule
\end{tabular}
\end{center}

\paragraph{Summary.} Cold-start NM-ROM on SIREN, Poisson-2D, N=128 now
reaches CP-decoder-class accuracy: median rel-L${}^2 = 2.88\times
10^{-3}$, $94\%$ of test parameters below $10^{-2}$, only $\sim 6\%$
catastrophic-tail outliers (high-frequency parameters where iter 2's
seed is itself poor). Compared to the starting point of rel-L${}^2 =
1.565$ at $z=0$, this is a $\boldsymbol{500\times}$ improvement in
median.

\paragraph{Key learnings.}
\begin{enumerate}
\item The cold-start failure is a \emph{basin-of-attraction} issue,
      not a Jacobian-rank or manifold-quality issue. The fix is at
      training time: anchor decode$(0)$ to the mean field and compress
      the encoder's latent distribution to be near 0.
\item Architectural changes (width, ROM-aware loss) shift the basin
      geometry \emph{independently} of manifold quality. They broke
      cold-start while improving warm-start.
\item The two-checkpoint composite solver gets the best of both:
      iter 2's reliable basin geometry seeds iter 5's
      ROM-quality manifold.
\end{enumerate}

\paragraph{Compute used.} Diagnostic: 20\,s. Iter 2: 10\,min A100.
Iter 3 (no retrain): 22\,min. Iter 4 (XL): 28\,min. Iter 5: 18\,min.
Iter 6: 17\,min. Iter 7 (no retrain): 7\,min. Phase 5 total
$\sim 100$\,min A100 time.

\end{document}
